% Source: https://artofproblemsolving.com/community/c3h1873894_what_are_odd_and_even_functions
% Stable identity: aops-topic-1873894-q01
Geometrically, $f$ is even when its graph is symmetric about the $y$-axis, and odd when its graph is symmetric about the origin (a half-turn about $O$ maps it to itself).

\textbf{(a)} $p(-x)=(-x)^4-3(-x)^2=x^4-3x^2=p(x)$, so $p$ is \emph{even}. $q(-x)=-x^3-x=-q(x)$, so $q$ is \emph{odd}. For $r$: $r(1)=2$ and $r(-1)=0$, so $r(-1)\ne r(1)$ and $r(-1)\ne-r(1)$; $r$ is \emph{neither}.

\textbf{(b)} If $f$ is even and odd, then $f(x)=f(-x)=-f(x)$, so $2f(x)=0$ and $f(x)=0$ for every $x$.

\textbf{(c)} \emph{Existence.} Put
\[
e(x)=\dfrac{f(x)+f(-x)}{2},\qquad o(x)=\dfrac{f(x)-f(-x)}{2}.
\]
Then $e(-x)=\dfrac{f(-x)+f(x)}{2}=e(x)$, $o(-x)=\dfrac{f(-x)-f(x)}{2}=-o(x)$, and $e(x)+o(x)=f(x)$.

\emph{Uniqueness.} If $f=e_1+o_1=e_2+o_2$, then $e_1-e_2=o_2-o_1$ is both even and odd, so by (b) it is $0$; hence $e_1=e_2$ and $o_1=o_2$.

For $r(x)=x^2+x$: $e(x)=\dfrac{(x^2+x)+(x^2-x)}{2}=x^2$ and $o(x)=\dfrac{(x^2+x)-(x^2-x)}{2}=x$. So $\boxed{r(x)=x^2+x\text{ with even part }x^2\text{ and odd part }x}$.
